\section{Law of Social Dynamics}

Auguste Comte proposed the Law of Social Dynamics\footnote{\url{https://gornahoor.net/library/AugusteComte.pdf}} where humanity passes necessarily through the three states of theological affirmation, metaphysical critique, and positive science or religion. In point of fact, from a historical perspective, not every society has passed through those three stages; furthermore, Comte has not demonstrated that every society is even capable of doing so. Nevertheless, it is helpful in organizing three different ways of understanding the world and ties in with the discussion of Instinct, Consciousness, and Science.

\paragraph{Theological Stage}


At this stage, society attributes everything to the Will of personified deities. At the level of Instinct this is perfectly rational. We moderns, who are accustomed to the motions of mechanical objects and to the attribution of motion to the hidden forces of gravitation and electro-magnetism, may not appreciate that for pre-modern peoples, motion is the result of conscious agents. In many respects, this is more rational than those moderns who tend to see things as always the result of impersonal forces, such as evolution, either biological or sociological.

Nowadays, we still see this attitude in the proliferation of conspiracy theories, which want to understand events in terms of deliberately planned actions. Again, the thrust behind such theories is understandable, since what transpires indeed is the result of the motives of conscious actors; it is just not as monolithic and perfectly planned as the conspiracy theorists would hold. Nevertheless, it suits the purposes of the agents of change to proscribe any discussion as to their identities and motives.

This stage may be divided into three sub-stages:

\subparagraph{Animism}


Animism or fetishism, attributes a will to everyday objects. This must be a transitional stage, since its explanatory power is meager and it leads to superstitions.

\subparagraph{Polytheism}


Polytheism attributes a will to many gods. Instinctually, this arises from the recognition and experience of transcendental or spiritual powers. For example, wind is understood as the cause of the motion of plants or inanimate objects, so fetishism can be dispensed with; instead, there is a god of the wind that supplies the explanation for movement. Similarly, the various powers that move the world or the mind of man, can be traced to transcendental experiences that are attributed to the gods.

\subparagraph{Monotheism}


Monotheism attributes a will to one god. Instinctually, this is the recognition of the ultimate unity of the world, despite its apparent multifariousness. There is trend in our day, mistakenly claiming to be based on Tradition, particularly Evola, to prefer polytheism to monotheism. Unfortunately, too, it justifies itself on utilitarian grounds rather than on the search for truth. This exposes the modernist roots of the idea, and Evola himself opposed such neo-paganism.

In point of fact, polytheism and monotheism, properly understood, go together. Even from a primitive point of view, there was always a ``chief'' god, who provided the ultimate unity. ``Zeus'', or Deus, simple means God. Jupiter, or ju-pater (zeus-pater), means Father-God. The reflective elements of pagan times, such as Plato, Aristotle, the Stoics, and Neo-Platonists, all accepted the ultimate unity of the cosmos under one God. The same can be said for the polytheism of the Hindus, with the Vedanta playing the other role.

Even Catholicism can be said to combine both; Comte called it the ``polytheism of the Middle Ages''. Keep in mind that, for the Ancients, a god was synonymous with immortal. Thus the Theotokos and the celestial hierarchy of medieval Catholicism would be considered gods by the Ancients.

\paragraph{Ideological Stage}


Comte calls this the metaphysical stage; however, since, in Tradition, metaphysics has a meaning distinct from that given it by secular philosophers, we prefer to call it the ideological stage, which is a more exact designation. This stage has the nature of critique and abstractions rather than offering a positive content. The beings that account for the world at the first stage are replaced with abstract principles, such a matter, progress, determinism, and so on. The theories or world views based upon such principles are characterized by their unfalsifiability, that is, no empirical fact can disprove them. This, along with their comprehensiveness that allegedly explains everything, make them very attractive. Someone in the grip of an ideological theory is very hard to dissuade, short of a ``conversion'' experience. Unfortunately, they often jump from one ideology to another.

\paragraph{Positive stage}


At this point, ideological explanations are dropped in favour of scientific explanations, based on observations of facts, experimentation, and the discovery of laws based on them. Comte sought to organize scientific knowledge in a hierarchy, based on their objects and a methodology appropriate to them. In particular, he developed the science of sociology and he is often called its ``Father''. Unfortunately, Comte is not read in any contemporary sociology department, but the ``other'' sociologist, Karl Marx, dominates. Marx and his legion of followers are reverts to the ideological stage and their revolutionary critique exerts a constant pressure to overthrow the natural order as intuitively understood at the theological stage.

Comte, on the other hand, eschews ideology as beyond the purview of positive or scientific laws. Instead, he seeks to ground the natural order on a firmer basis, so positivism does not undermine social stability. He writes:

\begin{quotex}
The best evidence of having attainted complete emancipation will be the rendering full justice to the past in all its phases. ... The surest sign of superiority, whether in persons or systems, is fair appreciation of opponents. And this must always be the tendency of social sciences when rightly understood, since the provision of the future is avowedly based upon systematic examination of the past. It is the only way in which the free and yet universal adoption of general principles of social reconstruction can ever be possible.
\end{quotex}


This makes positivism a reactionary system. In contrast, the revolutionary always looks to an imaginary future, when the laws of human nature will be somehow abrogated. For the revolutionary, Progress always means the overthrow of order. Yet, the fundamental motto of Positivism is \emph{Order and Progress}, but, for Comte, progress is simply the development of Order. ``The natural order necessarily contains within itself the germ of all possible progress.''

After Comte, positivism became ideological and adopted atheism, materialism, and determinism as unwarranted additions. Comte, however, moved in a different direction. He recognized that ethics, too, needs to be scientific, if sociology is to lead to progress. Furthermore, he created a religious system, strongly reminiscent of Catholicism, but without the supernaturalism, instead encompassing more of the greatest thinkers of history. For this, he was greatly criticized; however, it seems a reasonable development based on history.

\paragraph{A Further Post-Positive Stage}


Comte took the observation of outer phenomena as far as he could. However, it is an arbitrary decision to leave out inner experiences. Thus, Julius Evola could call for a ``positivism of metaphysics'' based on the study of inner experiences of a higher nature.


\flrightit{Posted on 2010-06-14 by Cologero}

